\begin{theorem}[Tame quadratic preparation]
\label{mod:cases-tamequadraticpreparation}
Let (K/F) be a prime-cyclic extension of nonarchimedean local fields, and
let (P:\texttt{PrimeCyclicPreparation}(F,K)) be the data supplied by the
ramified branch: its lower break, residue degree one, and an integral
monogenic uniformizer.  Assume that (K/F) is tamely ramified and that

\[
 [K:F]=2.
\]

Let \(Delta_F\) and \(Delta_K\) be local-constant functions satisfying
\texttt{IsDeltaFiniteLocalConstant}, and let \(chi_F\) and \(psi_F\) be
arbitrary exact multiplicative- and additive-conductor packages.  Then Lean
proves the literal fixed-character First Main identity

\[
 \Delta_K(\chi_F\circ N_{K/F},
          \psi_F\circ\operatorname{Tr}_{K/F})
 \prod_{\mu\in S(K/F)}\Delta_F(\mu,\psi_F)
 =
 \prod_{\mu\in S(K/F)}\Delta_F(\mu\chi_F,\psi_F),
\]

where (S(K/F)=\texttt{NormCharacter}(F,K)) is made finite from (P), and
both products include its identity character.  Equivalently, the conclusion
is exactly

\[
 \texttt{FirstMainIdentity}\ F\ K\ \Delta_F\ \Delta_K\
   \chi_F.\texttt{character}\ \psi_F.\texttt{character}.
\]

The proof first specializes the prepared break to zero using tameness.
Tameness, the prepared equality (e(K/F)=[K:F]), and degree two imply that
the residue characteristic is not two.  It then selects the
least-conductor member \(chi_{\min}\) of the complete norm-character orbit
of the arbitrary input \(chi_F\).  Lean constructs the exact norm pullback
of \(chi_{\min}\) using the canonical least multiplicative conductor, and
constructs the exact trace pullback of \(psi_F\) with conductor

\[
 n_K=[K:F]n_F+([K:F]-1).
\]

These data, together with the actual-conductor packages for every
\(mu\in S(K/F)\) and \(mu\chi_{\min}\), form
\texttt{FirstMainComputationalData}.  The already-completed theorem
\texttt{firstMain\_tameQuadratic} supplies its exact upper, norm-character,
and twist denominators and proves the full \texttt{deltaFinite} product.
The generic bridge
\texttt{firstMainIdentity\_of\_deltaFinite\_product} converts that product
to \texttt{FirstMainIdentity} for \(chi_{\min}\).

Finally, membership of \(chi_{\min}\) in the original norm-character orbit
and the proved invariance of both complete First Main sides under a fixed
norm-character twist transport the identity back to \(chi_F\).  All
conductor-zero, conductor-one, higher even, and higher odd calculations---in
particular every residual Hasse and quadratic phase---remain internal to
\texttt{Cases/TameQuadratic.lean} and are not reconstructed here.
\LeanFile{LanglandsFirstMainLemma/Cases/TameQuadraticPreparation.lean}
\lean{LanglandsFirstMainLemma.firstMain_tameQuadratic_of_isDeltaFinite}
\leanok
\SourceText{Propositions prop:tame-conductor-one-comparison and prop:tame-two-new; Section sec:first-main-completion.}
\uses{mod:ramification-primecyclicpreparation,mod:parameters-minimalorbitstationary,mod:cases-tamequadratic,mod:delta-errorterm}
\FormalizationContract
The public theorem takes the prepared ramified branch package directly; it
does not ask the caller for an already-minimal character or for norm/trace
pullback packages.  Its conclusion is the literal \texttt{FirstMainIdentity}
for the arbitrary original character.  Exact conductor packages and the
identity factor in the full norm-character products are retained.  This is
dependent routing around the completed case theorem, not a second quadratic
calculation.  Expected implementation size: 250--600 Lean lines.
\end{theorem}
